% =====================================================================
% TABLA COMPARATIVA  ---  COMPARTIDA (cada quien llena SUS filas)
% Rubrica: obligatoria. Columnas fijas segun el enunciado.
% FORMATO ACORDADO: metricas en % con un decimal; datasets abreviados.
% Agregar filas AL FINAL del bloque del eje propio.
% =====================================================================
\subsection{Tabla Comparativa}
\label{sec:tabla_comparativa}

The table summarizes the main experimental studies reviewed
across the three axes, including their modeling techniques, datasets or
environments, primary evaluation metrics, reported results, and identified
limitations. The comparison shows substantial variation in both prediction
tasks and experimental protocols. Horse-racing studies generally use
different jurisdictions, dataset sizes, target definitions, and evaluation
procedures, while studies from adjacent domains provide evidence for
alternative modeling techniques and problem formulations. Therefore, it
serves primarily to organize the evidence used to identify the research gaps
and methodological considerations discussed in the following sections, rather
than to establish a single ranking of the reviewed models. to see more go to
TABLE ~\ref{tab:comparativa}

\begin{table*}[!t]
\centering
\caption{Comparison of Related Work}
\label{tab:comparativa}
\scriptsize
\renewcommand{\arraystretch}{1.25}
\begin{tabular}{@{}p{0.9cm} p{0.8cm} p{2.5cm} p{2.7cm} p{1.8cm} p{1.5cm} p{3.8cm}@{}}
\toprule
\textbf{Ref.} & \textbf{Year} & \textbf{Technique} & \textbf{Dataset / environment} &
\textbf{Primary metric} & \textbf{Result} & \textbf{Identified limitation} \\
\midrule

\multicolumn{7}{@{}l}{\textit{Axis A --- domain (Int. 1)}} \\
\midrule

\cite{ejeA_davoodi2010} & 2010 & ANN (5 training algorithms) &
Horse racing, Aqueduct NY (100 races) &
Accuracy & 77.0\,\% &
Single track; one month of data \\

\cite{ejeA_schumaker2018} & 2018 & SVR &
Harness racing, Northfield Park (Ohio) &
Accuracy & 50.4\,\% &
Single harness-racing track; performance depends on race-history window \\

\cite{ejeA_gupta2023} & 2023 & R F + SMOTE &
Horse racing, Turf Club of India (14,750 records) &
Accuracy / F1-score / ROC-AUC &
93.1\,\% / 92.9\,\% / 97.6\,\% &
Single site; different jurisdiction from ours \\

\cite{ejeA_li2025} & 2025 & Linear / log-linear regression &
Horse racing, 3 UK racecourses &
$R^2$ & 0.995 &
Models winning-time consistency rather than the race winner \\

\midrule
\multicolumn{7}{@{}l}{\textit{Axis B --- techniques (Int. 2)}} \\
\midrule

\cite{ejeB_kholkine2020} & 2020 & XGBoost + SHAP &
Road cycling, Tour of Flanders &
Top-10 hits / 10 races & 6/10 (2018) &
Small dataset; several performance across race editions \\

\cite{ejeB_peng2020} & 2020 & RankNet (LSTM encoder-decoder) vs. RF/SVM/XGBoost &
Car racing, IndyCar (Indy500) &
Top-1 accuracy / MAE &
79.0\,\% / 0.94 &
Classical ML baselines did not consistently beat the no-change baseline \\

\cite{ejeB_dumitrescu2022} & 2022 & Penalised Logistic Tree Regression (PLTR) &
Credit scoring (4 datasets) &
AUC & 85.7\,\% &
Cross-domain; several model performance across datasets \\

\cite{ejeB_vuko2025} & 2025 &
Logistic regression / Lasso / Ridge / Elastic Net&
Going-concern prediction, Croatia (891 obs.) &
AUC-ROC & 84.2\,\% &
Small, imbalanced dataset; differences between models not statistically significant \\

\cite{ejeB_krzyszton2026} & 2026 & SVM + RF + GB ensemble &
Formula 1 race results (OpenF1, 1,635 records) &
Macro F1-score & 77.8\,\% &
Winner class hardest; only 2023--2024 seasons \\

\cite{ejeB_silverman2013} & 2013 & Conditional logistic regression + frailty (GPU) &
Horse racing, 3,681 races &
GPU/OpenMP speedup & by a time factor of 2.68 &
No accuracy or AUC reported; computational result only \\

\midrule
\multicolumn{7}{@{}l}{\textit{Axis C --- dataset (Int. 3)}} \\
\midrule

\cite{ejeA_gupta2023} & 2012 &
Hierarchical Bayesian model (running speed) &
Horse racing, Hong Kong (2,973 races) &
Accuracy (fastest horse) & 21.6\,\% &
Predicts fastest horse, not official winner; period unspecified \\

\cite{ejeC_chavan2024} & 2024 &
Logistic Regression / Decision Tree / RF / SVM / Dummy Classifier &
Horse racing, 6,348 races / 79,447 records &
Accuracy / F1-score &
92.0\,\% / 51.0\,\% &
Random 80/20 split; Dummy Classifier also achieves 92.0\,\% accuracy \\

\bottomrule
\end{tabular}
\end{table*}

The comparison table uses several standard machine-learning abbreviation
and terminology. \textit{Artificial Neural Network} (ANN) refers to a
neural-network model for learning nonlinear relationships between input
features and a target. \textit{Support Vector Regression} (SVR) is a
support-vector method for predicting continuous outcomes, while
\textit{Support Vector Machine} (SVM) is its commonly used classification
counterpart. \textit{Random Forest} (RF) is an ensemble of decision trees
whose predictions are aggregated, and \textit{Gradient Boosting} (GB)
sequentially builds trees to correct previous prediction errors.
\textit{XGBoost} is an optimized gradient-boosting implementation.

\textit{SMOTE} (Synthetic Minority Over-sampling Technique) is a
class-imbalance preprocessing method that generates synthetic minority-class
examples. Other terms in the table describe methodological combinations rather
than individual predictive models. For example, \textit{SHAP} (SHapley Additive
exPlanations) is an interpretability method used to estimate the contribution 
of individual features to model predictions. These terms describe either
predictive models or supporting techniques and should not be interpreted as
equivalent types of algorithms.

% Formato de fila (copiar y rellenar):
% \cite{clave} & 2022 & Tecnica & Dataset & Metrica & 27.4\,\% & Limitacion \\

% NOTA: los articulos teoricos o de revision que no reportan resultado
% experimental NO van en la tabla; se discuten en el texto y se explica
% en una linea por que quedan fuera.